Модель включает пять сущностей: \textbf{Area}, \textbf{Zone},
\textbf{Camera}, \textbf{Event} и \textbf{CameraTelemetry}.

\subsubsection{Area}

\textbf{Area} --- охраняемая территория, на которой развернута система
видеонаблюдения.

Основные атрибуты территории:

\begin{itemize}
  \item \emph{area\_code} --- предметный идентификатор территории;
  \item \texttt{name} --- название территории;
  \item \texttt{area\_type} --- тип территории;
  \item \texttt{address} --- адрес или описание местоположения;
  \item \texttt{description} --- дополнительное описание территории.
\end{itemize}

Одна территория может содержать несколько зон наблюдения.

\subsubsection{Zone}

\textbf{Zone} --- часть охраняемой территории, внутри которой установлены
камеры и фиксируются события.

Основные атрибуты зоны:

\begin{itemize}
  \item \emph{zone\_code} --- идентификатор зоны внутри территории;
  \item \emph{area\_code} --- идентификатор территории, к которой относится
        зона;
  \item \texttt{name} --- название зоны;
  \item \texttt{zone\_type} --- тип зоны;
  \item \texttt{importance\_level} --- уровень важности зоны;
  \item \texttt{description} --- дополнительное описание зоны.
\end{itemize}

Идентификатор зоны является составным: зона именуется кодом территории и
кодом зоны внутри этой территории. Такой вариант соответствует
идентификационно-зависимой сущности: зона не рассматривается отдельно от
территории, к которой она относится.

\subsubsection{Camera}

\textbf{Camera} --- камера видеонаблюдения, установленная в зоне.

Основные атрибуты камеры:

\begin{itemize}
  \item \emph{serial\_number} --- предметный идентификатор камеры;
  \item \texttt{name} --- название камеры;
  \item \texttt{model} --- модель камеры;
  \item \texttt{ip\_address} --- сетевой адрес камеры;
  \item \texttt{status} --- текущее состояние камеры;
  \item \texttt{position} --- положение и ориентация камеры внутри зоны;
  \item \texttt{settings} --- настройки видеопотока и аналитики.
\end{itemize}

Камера относится к одной зоне наблюдения. При этом одна зона может содержать
несколько камер. Внутренние зоны детекции камеры, используемые алгоритмами
аналитики, не являются зонами предметной области и входят в состав
атрибута \texttt{settings}.

\subsubsection{Event}

\textbf{Event} --- аналитическое событие, возникшее в зоне наблюдения.

Основные атрибуты события:

\begin{itemize}
  \item \emph{event\_number} --- номер события внутри зоны;
  \item \emph{zone\_code} --- идентификатор зоны, в которой произошло
        событие;
  \item \emph{area\_code} --- идентификатор территории события;
  \item \texttt{occurred\_at} --- время возникновения события;
  \item \texttt{event\_type} --- тип события;
  \item \texttt{severity} --- уровень значимости события;
  \item \texttt{confidence} --- уверенность аналитического алгоритма;
  \item \texttt{payload} --- составные данные события.
\end{itemize}

Идентификатор события является составным: номер события уникален внутри зоны,
а зона определяется кодом зоны и кодом территории. Событие может быть
зафиксировано одной или несколькими камерами. Например, один факт пересечения
линии может подтверждаться двумя камерами, установленными в одной зоне.

Атрибут \texttt{payload} содержит параметры, зависящие от типа события:
параметры движения, сведения о пересеченной линии, информацию о потере
сигнала или массив обнаруженных объектов. В ER-модели это составной
вариативный атрибут. В разных физических моделях он реализуется по-разному:
как \texttt{jsonb}, как набор нормализованных таблиц или как вложенный
документ.

\subsubsection{CameraTelemetry}

\textbf{CameraTelemetry} --- регулярная запись технического состояния камеры.
Эта сущность хранит историю измерений, а не только текущее состояние камеры.

Основные атрибуты телеметрии:

\begin{itemize}
  \item \emph{serial\_number} --- идентификатор камеры;
  \item \emph{recorded\_at} --- время фиксации телеметрии;
  \item \texttt{status} --- состояние камеры на момент фиксации;
  \item \texttt{metrics} --- составные технические показатели камеры.
\end{itemize}

Идентификатор записи телеметрии состоит из идентификатора камеры и времени
измерения. Атрибут \texttt{metrics} включает технические показатели, например
температуру, загрузку процессора, использование памяти, битрейт, потери
пакетов, задержку и время непрерывной работы.

\subsubsection{Итоговый набор сущностей}

Итоговый набор сущностей ER-модели:

\begin{itemize}
  \item \textbf{Area} --- охраняемая территория;
  \item \textbf{Zone} --- зона наблюдения внутри территории;
  \item \textbf{Camera} --- камера видеонаблюдения;
  \item \textbf{Event} --- аналитическое событие;
  \item \textbf{CameraTelemetry} --- запись технического состояния камеры.
\end{itemize}
